\section{Phonetic Bootstrap Testbench}

\subsection{Purpose}

The project can now begin phonetic bootstrapping in a controlled sense. This does not mean that a sign has been read as a sound. It means that a reproducible testbench now assigns abstract variables to signs and extracts the inscriptions where future sound values, morphemes, titles, classifiers, or proper-name readings can be falsified.

The notation is deliberately conservative. A variable such as \texttt{C01}, \texttt{T01}, or \texttt{R02} is not a phoneme. It is a placeholder for a sign whose function is constrained enough that a phonetic hypothesis can later be tested against every occurrence.

\subsection{Method}

The script \texttt{scripts/phonetic-bootstrap-testbench.py} combines:

\begin{itemize}
  \item the 607 medium/high phonetic-readiness reconstructions from the structural model;
  \item functional proto-glosses for signs such as \texttt{820}, \texttt{861}, \texttt{817}, \texttt{740}, and \texttt{520};
  \item cross-frame filler reuse from the slot-paradigm model;
  \item minimal alternation pairs within the same structural frame.
\end{itemize}

Each sign receives a stable abstract variable. The prefix marks only the current functional class:

\begin{itemize}
  \item \texttt{C}: classifier/title-like sign;
  \item \texttt{T}: terminal title/suffix-like sign;
  \item \texttt{M}: formula marker;
  \item \texttt{A}: slot modifier;
  \item \texttt{R}: possible root or unresolved sign.
\end{itemize}

\subsection{Corpus-Wide Result}

\begin{table}[htbp]
\centering
\begin{tabular}{lr}
\toprule
\textbf{Metric} & \textbf{Value} \\
\midrule
Medium/high ready reconstructions & 607 \\
Abstract sign variables & 338 \\
Distinct reading units & 627 \\
Strong onomastic anchors & 7 \\
Cross-frame phonetic anchors & 31 \\
Affix/formula anchors & 75 \\
Prime-entity fillers & 85 \\
Abstract reconstructions exported & 607 \\
Minimal alternation tests & 300 \\
\bottomrule
\end{tabular}
\caption{Abstract phonetic bootstrap summary.}
\label{tab:phonetic-bootstrap-summary}
\end{table}

\subsection{Variable Map}

The strongest currently constrained variables are shown in Table~\ref{tab:phonetic-variable-map}. These are the variables that future proposed readings must explain first.

\begin{table}[htbp]
\centering
\begin{tabular}{lll}
\toprule
\textbf{Sign} & \textbf{Variable} & \textbf{Current functional class} \\
\midrule
820 & \texttt{C01} & classifier/title-like \\
861 & \texttt{C02} & classifier/title-like \\
817 & \texttt{C03} & classifier/title-like \\
740 & \texttt{T01} & terminal title/suffix-like \\
520 & \texttt{T02} & terminal title/suffix-like \\
235 & \texttt{M01} & formula marker / slot modifier \\
390 & \texttt{M02} & formula marker / slot modifier \\
002 & \texttt{A01} & slot modifier \\
240 & \texttt{A02} & slot modifier \\
176 & \texttt{R02} & possible root or unresolved sign \\
\bottomrule
\end{tabular}
\caption{Selected abstract variables. These are not sound values.}
\label{tab:phonetic-variable-map}
\end{table}

\subsection{Strong Onomastic Anchors}

The most important new class is the strong onomastic anchor: a prime-entity filler that appears in a semantically narrow slot, occurs across more than one site, and participates in minimal alternation evidence. These are the closest analogues, in our current data, to the proper-name strategy used in successful decipherments.

\begin{table}[htbp]
\centering
\footnotesize
\begin{tabular}{llllr}
\toprule
\textbf{Unit} & \textbf{Abstract} & \textbf{Role} & \textbf{Frame evidence} & \textbf{Score} \\
\midrule
176 & \texttt{R02} & prime entity filler & \texttt{002|740}; \texttt{<START>|741} & 0.790 \\
032-220 & \texttt{A07-A09} & prime entity filler & \texttt{002|740}; \texttt{<START>|233} & 0.757 \\
590-390 & \texttt{A08-M02} & prime entity filler & \texttt{002|740} & 0.737 \\
240-482 & \texttt{A02-R16} & prime entity filler & \texttt{002|740} & 0.664 \\
235-222 & \texttt{M01-R31} & prime entity filler & \texttt{002|740} & 0.623 \\
240-176 & \texttt{A02-R02} & prime entity filler & \texttt{002|740} & 0.608 \\
235-840-032 & \texttt{M01-R10-A07} & prime entity filler & \texttt{002|740} & 0.600 \\
\bottomrule
\end{tabular}
\caption{First strong onomastic-style anchors for controlled phonetic testing.}
\label{tab:strong-onomastic-anchors}
\end{table}

\subsection{Example Abstract Reconstructions}

The output file \texttt{outputs/abstract\_phonetic\_reconstructions.csv} renders 607 inscriptions in abstract variables. A few examples:

\begin{itemize}
  \item \textbf{L-41, Lothal}:\\
  {\footnotesize\texttt{817-002-233-740-090} to \texttt{C03-A01-A04-T01-A10}}.\\
  The sign \texttt{233} is the prime entity filler inside \texttt{002 X 740}.
  \item \textbf{M-708, Mohenjo-daro}:\\
  {\footnotesize\texttt{140-390-002-061-740} to \texttt{M07-M02-A01-R29-T01}}.\\
  This also preserves the \texttt{002 X 740} entity slot.
  \item \textbf{M-1114, Mohenjo-daro}:\\
  {\footnotesize\texttt{920-060-741-002-140-740} to \texttt{M04-A06-A03-A01-M07-T01}}.\\
  This is a longer prime-entity seal formula.
\end{itemize}

This is the first point where the project can speak of sentence-like reconstruction with a phonetic layer: the formula is structurally and semantically constrained, while the sound layer remains abstract.

\subsection{Minimal Alternation Protocol}

The file \texttt{outputs/phonetic\_minimal\_tests.csv} contains the top 300 same-frame alternations. These are essential because a proposed reading is not allowed to explain only one attractive inscription. It must survive the alternation set.

For example, if two fillers occupy the same \texttt{002|740} prime-entity slot and differ by one sign, then a proposed phonetic value for that sign must preserve the same frame. The alternation may indicate a phonemic contrast, an affix, a classifier, an allograph, a numeral, or a graphic modifier. The testbench does not decide which one is correct; it gives us the falsification grid.

\subsection{Interpretation}

The current working reconstruction pipeline is now:

\begin{center}
\texttt{STRUCTURAL FRAME -> SEMANTIC OPTIONS -> ABSTRACT PHONETIC VARIABLES -> MINIMAL TESTS}
\end{center}

The next empirical breakthrough will probably not come from guessing a sound for a sign in isolation. It should come from taking a strong anchor such as \texttt{176}, \texttt{590-390}, or \texttt{235-222}, collecting every artifact and image context where it appears, and asking whether one proposed semantic or phonetic value predicts the whole set.

\subsection{Outputs}

\begin{itemize}
  \item \texttt{outputs/phonetic\_bootstrap\_summary.csv}
  \item \texttt{outputs/phonetic\_variable\_map.csv}
  \item \texttt{outputs/phonetic\_reading\_units.csv}
  \item \texttt{outputs/abstract\_phonetic\_reconstructions.csv}
  \item \texttt{outputs/phonetic\_minimal\_tests.csv}
  \item \texttt{outputs/phonetic\_bootstrap\_testbench.tex}
\end{itemize}
